\subsection{Обзор аналогов}
\label{sec:analysis:analogues}

Следует отметить, что существует множество информационных систем, поддерживающих удаленное взаимодействие в менеджменте студий \textit{3D}-печати и контакт между участниками процесса (сотрудниками и клиентами). Некоторые системы акцентируют внимание на простоте загрузки моделей и автоматическом расчете стоимости (\textit{Instant Quote}), в то время как другие предлагают более сложные и детализированные интерфейсы для управления парком оборудования (\textit{print farm}), складскими остатками филамента и интеграции с \textit{CAD}-системами.

Исследование существующих аналогичных систем является важной частью процесса проектирования информационной системы для студии \textit{3D}-печати. Учет преимуществ и недостатков всех аналогов позволит выявить среднее значение по функциональным элементам и получить оптимальное решение в рамках курсового проекта.

В рамках разработки АИС студии \textit{3D}-печати планируется выявление оптимального функционала, необходимого для следующих элементов системы:
\begin{itemize}
    \item домашняя страница администратора с панелью управления активными принтерами;
    \item личный кабинет клиента с историей заказов и статусами выполнения;
    \item списки хранимых данных о пользователях, материалах (\textit{filaments}) и типах пластика;
    \item общий каталог доступных материалов с актуальными ценами;
    \item справочный каталог принтеров с их текущим статусом (\textit{active state});
    \item функционал для сотрудников по списанию материала и смене статусов заказов;
    \item аналитическая панель (\textit{dashboard}) для роли Аналитика, отображающая статистику и графики.
\end{itemize}
Оптимальность полученного функционала позволит сократить расходы на разработку и внедрение.

\textbf{1.2.1}\hphantom{22}Крупные агрегаторы и студии, такие как \textit{Studia3D}, часто предлагают широкий выбор услуг, включая не только печать, но и 3D-сканирование, моделирование и продажу оборудования. Их онлайн-порталы используют комплексный подход для распространения информации.
На рисунке \ref{fig:studia3d_homepage} представлена главная страница портала \textit{Studia3D}.
 
\begin{figure}[H]
    \centering
    \includegraphics[width=1\textwidth]{studia3d_homepage.png}
    \caption{Сайт компании \textit{Studia3D}}
    \label{fig:studia3d_homepage}
\end{figure}

При изучении такого портала можно выявить ряд недостатков в контексте автоматизации конкретной студии. Интерфейс перегружен информационными блоками: новостями индустрии, рекламой принтеров для продажи и множеством калькуляторов. Это может отклонять внимание от стандартного \textit{CJM (Customer Journey Map)} пользователя, который просто хочет распечатать деталь.

Отсутствие строгого фокуса на процессе «заказ – печать», характерное для маркетплейсов, может порождать ряд проблем для локальной студии. Избыточные формы обратной связи, всплывающие окна консультантов и смешивание услуг печати с интернет-магазином товаров могут способствовать снижению качества пользовательского опыта (\textit{UX}) для постоянного клиента, которому важна скорость оформления.

Сайт пользуется популярностью в регионе за счет бренда и энциклопедического контента. Однако сложность интерфейса и обилие лишней информации следует учесть при проектировании более грамотного и лаконичного интерфейса собственной АИС.

\textbf{1.2.2}\hphantom{22}Сайт сервиса \textit{CanTouch}, специализирующегося на промышленной и частной \textit{3D}-печати, акцентирует внимание на простоте онлайн-расчета. Интерфейс сайта приведен на рисунке \ref{fig:cantouch_homepage}.

\begin{figure}[H]
    \centering
    \includegraphics[width=1\textwidth]{cantouch_homepage.png}
    \caption{Сайт компании \textit{CanTouch}}
    \label{fig:cantouch_homepage}
\end{figure}

Описываемый сайт является противоположностью перегруженным порталам. Лаконичность интерфейса, клиентоориентированность, «чистый экран» и фокус на целевом действии улучшают пользовательский опыт. Система сразу предлагает выбрать технологию и материал, отсекая лишние шаги, а затем сделать заказ.

Стоит учитывать предоставляемое сайтом архитектурное решение по выбору материалов. На рисунке \ref{fig:cantouch_order} приведен экран заказа печати.

Пользователь видит не список материалов, а форму заполнения для обратной связи, для взаимодействия Клиент-Сотрудник.

Отметим, что на сайте отсутствует регистрация, поэтому, вероятнее всего, ни о какой глубокой автоматизации не может идти речи, а данные хранятся, как и в большинстве подобных студий, в таблицах \textit{Excel}.

Также видно, что потребитель сам должен загрузить на какое-то стороннее хранилище файл модели и указать ссылку на него. Это усложняет процесс заказа, но позволяет избежать хранения больших файлов на стороне сервиса.

\begin{figure}[H]
    \centering
    \includegraphics[width=1\textwidth]{cantouch_order.png}
    \caption{Экран заказа на сервисе \textit{CanTouch}}
    \label{fig:cantouch_order}
\end{figure}

Однако данный сервис ориентирован исключительно на клиента («фронтенд»). В нем скрыта внутренняя кухня: клиент не видит, какой конкретно принтер будет печатать его деталь, и не может отследить историю списания материалов. Для целей данного курсового проекта, где требуется управление сотрудниками, оборудованием и аналитикой, функционала только «витрины» недостаточно.

\textbf{1.2.3}\hphantom{22}Сайт-производитель профессионального управляющего ПО, предлагающий комплексное \textit{B2B}-решение для управления 3D-фермами и производством, предоставляется компанией \textit{DigiFabster}. Интерфейс системы приведен на рисунке \ref{fig:digi_homepage}.

Данная система соответствует современным стандартам \textit{CRM} и \textit{ERP} для аддитивного производства. Имеются вкладки под аналитику продаж, управление очередью печати (\textit{Order Queue}), складской учет филамента и мониторинг принтеров.

При сравнении с предыдущими сайтами, решение вынесения производственного учета и аналитики в отдельный интерфейс администратора является стратегическим ходом. В ходе поиска нужного заказа в массиве данных, обоснованном подходом со стороны функционала (фильтры по статусам, материалам), а не простого списка, повышается управляемость студией.

\begin{figure}[H]
    \centering
    \includegraphics[width=1\textwidth]{digi_homepage.png}
    \caption{Сайт компании \textit{DigiFabster}}
    \label{fig:digi_homepage}
\end{figure}

Система предлагает модули для контроля загрузки оборудования и формирования отчетности. Таким образом, уровень включенности сотрудника (\textit{Employee}) в процесс функционирования системы повышается на несколько порядков. Важно отметить, что такая система предлагает широкий ассортимент управленческих деталей, необходимых для принятия решений:

\begin{itemize}
    \item автоматическое обновление статусов (\textit{Order Status});
    \item расчет стоимости на основе геометрии модели;
    \item генерация финансовых отчетов и графиков эффективности.
\end{itemize}

Именно наличие мощного аналитического блока делает подобные системы ценными для бизнеса, подтверждая необходимость выделения роли Аналитика (\textit{Analyst}) в проектируемой системе.

Недостатком данной системы для малого бизнеса является её высокая стоимость и избыточность функционала (например, интеграция с почтовыми службами доставки США), что делает её внедрение в небольшой студии нецелесообразным. Однако архитектурные принципы разделения ролей и учета ресурсов, заложенные в \textit{DigiFabster}, будут взяты за основу при проектировании серверной части (\textit{backend}) нашей АИС.

На основе проведенного анализа можно сделать вывод, что идеальная АИС для студии \textit{3D}-печати должна сочетать в себе:
\begin{itemize}
    \item \textbf{Простоту клиентского интерфейса} (по аналогии с \textit{CanTouch}), позволяющую пользователю быстро оформить заказ и видеть его актуальный статус.
    \item \textbf{Мощный административный модуль} (по упрощенной модели \textit{DigiFabster}), который позволит сотрудникам управлять типами филамента (\textit{Filament Types}), вести учет остатков (\textit{Residue}) и контролировать историю действий (\textit{History}).
    \item \textbf{Инструменты бизнес-аналитики}, предоставляющие отдельной роли (\textit{Analyst}) доступ к агрегированным данным и графикам для принятия управленческих решений без риска вмешательства в производственный процесс.
    \item \textbf{Ролевую модель}, исключающую перегрузку интерфейса: каждый участник системы видит только те инструменты, которые необходимы для его задач.
\end{itemize}